\documentclass[11pt,fontset=fandol]{ctexart}
\usepackage[a4paper,margin=1in]{geometry}
\usepackage{amsmath,amsfonts,amssymb}
\usepackage{graphicx}
\usepackage{booktabs}
\usepackage{array}
\usepackage{caption}
\usepackage{enumitem}
\usepackage[hidelinks]{hyperref}
\usepackage{microtype}

\title{Lazy 环 + 单向 shortcut：图1同轴叠加与敏感性分析}
\author{}
\date{}

\setlist{nosep}
\setlength{\textfloatsep}{8pt plus 2pt minus 2pt}
\setlength{\floatsep}{6pt plus 2pt minus 2pt}
\setlength{\intextsep}{6pt plus 2pt minus 2pt}
\setlength{\abovecaptionskip}{4pt}
\setlength{\belowcaptionskip}{0pt}
\newcolumntype{L}[1]{>{\raggedright\arraybackslash}p{#1}}

\begin{document}
\maketitle

\section*{摘要}
本报告聚焦“图1同轴叠加”与“敏感性分析”，不重复原扩展报告的完整内容。
默认示例参数为 $N=100$，$K\in\{2,4\}$，$q=2/3$，$\beta=0.01$，$\rho=1.0$，模式为 lazy\_selfloop。
图1改为上下两幅（上 $K=2$、下 $K=4$），每幅将 $f(t)$ 与三段窗口（peak1/valley/peak2）中的四类轨迹组成以同一主轴显示；
并给出阈值、窗口位置/宽度、MC（Monte Carlo）不确定性的敏感性结论。

\section*{模型与符号速查}
\begin{itemize}
  \item $N$：环上节点数；$K=2/4$ 表示每步可向 $\pm r$ 移动，$r=\lfloor K/2 \rfloor$，并带 lazy 停留与单向 shortcut。
  \item jump-over：轨迹跨过吸收点（target）但未吸收的事件；与 shortcut 使用独立统计。
  \item 四类：$\texttt{C0J0}$、$\texttt{C1pJ0}$、$\texttt{C0J1p}$、$\texttt{C1pJ1p}$，分别表示
        “无 shortcut/无 jump-over”、“有 shortcut/无 jump-over”、“无 shortcut/有 jump-over”、“有 shortcut/有 jump-over”。
  \item $C1p/J1p$ 表示 “至少一次 shortcut/jump-over”，$C0/J0$ 表示 “从未发生”。
\end{itemize}

\section*{参数与窗口定义}
默认示例参数：$N=100$，$K\in\{2,4\}$，$q=2/3$，$\beta=0.01$，$\rho=1.0$，模式为 lazy\_selfloop。
其中 $q$ 为 lazy 留停概率，$\beta$ 控制 shortcut 强度，$\rho$ 为首达吸收率；$t$ 为离散时间步。

窗口宽度采用既有流程中的统一规则：给定前两峰 $t_1,t_2$ 与谷底 $t_v$，
\[
\Delta=\max\{1,\lfloor \delta_{\text{frac}}(t_2-t_1)\rfloor\},\quad \delta_{\text{frac}}=0.05,
\]
windows 定义为
$[t_1-\Delta,t_1+\Delta]$、$[t_v-\Delta,t_v+\Delta]$、$[t_2-\Delta,t_2+\Delta]$。
窗口中心 $(t_1,t_v,t_2)$ 由精确 $f(t)$ 的峰谷检测规则自动给出，因此不是随机选取，
实现细节记录在分析代码中。
本报告对 $K=2$ 与 $K=4$ 分别从各自的精确 $f(t)$ 检测得到 $(t_1,t_v,t_2)$，并用同一固定规则计算 $\Delta$；
上下两幅各用自身窗口，使栈叠条与该 $K$ 的峰谷结构对齐。
窗口选择理由分为定性与定量两部分：定性上，双峰分布对应快/慢通道，
peak1/peak2 抽取两端的“快/慢”动力学，valley 抽取两通道分离最强的时段，
对 MC 轨迹比例 $P(\text{class}\mid T\in\text{window})$ 最具解释力。
定量上，窗口中心来自峰谷检测、宽度与峰间距成比例，且在 $\Delta=\pm 1,2,5$ 的扰动下比例变化仍为 $10^{-3}$ 量级，
说明结论不依赖精确边界。
此外，表~\ref{tab:window-params} 给出三段窗口的中心、区间、宽度与窗口内样本量，
用于量化说明窗口位置与样本规模并非随意设定。

% Auto-generated by export_window_table.py
\begin{table}[!htbp]
\centering
\scriptsize
\setlength{\tabcolsep}{3pt}
\caption{窗口参数与样本量（$\beta=0.01, N=100$；窗口中心来自 $K=4$ 峰谷检测）。}
\label{tab:window-params}
\resizebox{\linewidth}{!}{%
\begin{tabular}{lrrrrrrrrr}
\toprule
window & $t_{\mathrm{center}}$ & $t_L$ & $t_R$ & $\Delta$ & width & $f_{K=2}(t_c)$ & $f_{K=4}(t_c)$ & $n_{K=2}$ & $n_{K=4}$ \\
\midrule
peak1 & 32 & 9 & 55 & 23 & 46 & 1.22e-04 & 1.10e-04 & 2337 & 924 \\
valley & 84 & 61 & 107 & 23 & 46 & 1.26e-04 & 9.52e-05 & 1981 & 940 \\
peak2 & 495 & 472 & 518 & 23 & 46 & 1.62e-04 & 6.23e-04 & 5796 & 5741 \\
\bottomrule
\end{tabular}}
\end{table}


\section*{复现说明（最短命令链）}
在本报告目录内运行：
\begin{verbatim}
make fig2
make sensitivity
make pdf
\end{verbatim}
输出路径：图1在 \texttt{figures/}，敏感性结果在 \texttt{outputs/sensitivity/}，LaTeX 构建在 \texttt{build/}。

\section*{数据来源表}
\begin{table}[!htbp]
\centering
\scriptsize
\setlength{\tabcolsep}{3pt}
\begin{tabular}{@{}L{0.18\linewidth} L{0.25\linewidth} L{0.25\linewidth} L{0.24\linewidth}@{}}
\toprule
图/表 & 生成程序 & 输入 & 输出/命令 \\
\midrule
图1（上下两幅） & \path{code/plot_fig2_overlap_binbars.py} & \path{data/fig2_bins_bars_beta0.01.json}, \path{data/ft_beta0.01.csv} & \path{figures/fig2_overlap_binbars_beta0.01_x1350.pdf}；\texttt{make fig2} \\
阈值敏感性 & \path{code/sensitivity_thresholds.py} & \path{data/ft_beta0.01.csv} & \path{outputs/sensitivity/threshold_sweep_heatmap.pdf} \\
窗口/宽度敏感性 & \path{code/sensitivity_bin_intervals.py} & \path{data/fig2_bins_bars_beta0.01.json} + \path{cond_by_t.csv} & \path{outputs/sensitivity/bin_shift_plot.pdf} \\
MC 不确定性 & \path{code/sensitivity_mc_uncertainty.py} & \path{data/fig2_bins_bars_beta0.01.json} + \path{cond_by_t.csv} & \path{outputs/sensitivity/mc_uncertainty_plot.pdf} \\
\bottomrule
\end{tabular}
\end{table}

\section{同轴叠加 + 窗口栈叠条}
图~\ref{fig:fig1-overlap-cn} 改为上下两幅：上图为 $K=2$，下图为 $K=4$。
每幅将首达时间分布 $f(t)$ 与窗口内轨迹组成同轴显示；
其中 $K=2$ 曲线为深蓝/黑、$K=4$ 曲线为红色（颜色一致），用于直观对比双峰/谷差异；
横轴 $t$ 为离散时间步（取值由 \path{data/ft_beta0.01.csv} 给定，本图为 $t=1\ldots 1350$ 以增强可读性）。
三段栈叠条对应 peak1/valley/peak2 三个窗口，它们的位置并非随机，
而是由各自 $f(t)$ 的两峰一谷检测得到的 $(t_1,t_v,t_2)$ 与固定规则宽度 $\Delta$ 构成的区间 $[t_L,t_R]$。
选择这三段窗口是为了直接对比“早峰/谷底/晚峰”内的轨迹组成：
peak1/peak2 对应快/慢通道的贡献，valley 反映两通道分离程度，
用窗口内的 MC 比例 $P(\text{class}\mid T\in\text{window})$ 能将三种动力学阶段分开观察。
窗口选择的合理性通过“窗口位置与宽度敏感性”定量检验：$\Delta=\pm 1,2,5$ 的扰动下，
比例变化仍在 $10^{-3}$ 量级，说明结论不依赖精确边界。具体窗口中心与区间见图~\ref{fig:fig1-overlap-cn} 标注。
与原报告的“橙/绿/蓝窗口底色”不同，此处用三段全高栈叠条替代 bin：
每个窗口对应一条栈叠 bar（宽度覆盖 $[t_L,t_R]$），内部用四种颜色表示
$P(\text{class}\mid T\in\text{window})$；上下两幅分别对应各自的 $K$ 与各自窗口。
若某类比例为 0，对应颜色不会出现。若某些时刻的条件概率被遮罩，
脚本将回退到窗口级汇总比例，并在 JSON meta 中记录来源。

\begin{figure}[!htbp]
  \centering
  \includegraphics[width=0.95\linewidth]{figures/fig2_overlap_binbars_beta0.01_x1350.pdf}
  \caption{$f(t)$ 与窗口内四类比例的同轴叠加（栈叠 bar 替代窗口底色），上下两幅分别为 $K=2$（深蓝/黑）与 $K=4$（红色）。
  $t$ 为离散时间步（横轴展示 $t=1\ldots 1350$ 以增强可读性）；标注的 $t_1,t_v,t_2$ 为各自 $K$ 的两峰一谷中心位置；
  三段栈叠条对应 peak1/valley/peak2 窗口，区间由各自 $(t_1,t_v,t_2)$ 检测并按固定规则宽度 $\Delta$ 构造,
  条内颜色为 $P(\text{class}\mid T\in\text{window})$（对应各自 $K$）。}
  \label{fig:fig1-overlap-cn}
\end{figure}

\subsection*{窗口栈叠条揭示的双峰机制}

\textbf{现象。} 图~\ref{fig:fig1-overlap-cn} 两幅都显示 $f(t)$ 在 $t_1$ 附近有早峰、
$t_v$ 附近有谷底、$t_2$ 附近有晚峰，且 $t_2 \gg t_1$（$K=2$ 时 $t_2/t_1\approx 21$,
$K=4$ 时 $\approx 15$）。双峰不是数值假象：每条窗口栈叠条上的类别组成
\emph{告诉我们哪类轨迹在贡献哪一个峰}。

\textbf{证据。} 在两幅中沿 peak1 $\to$ peak2 阅读各栈叠条,类别组成发生反转。
$K=2$ 那一幅 peak1 几乎是单一的 $\texttt{C1pJ0}$ 群体（至少用过一次 shortcut、
未跨过 target);peak2 几乎是单一的 $\texttt{C0J0}$ 群体（从未用 shortcut、
也未 jump-over）;valley 窗口正好位于这两种组成的过渡带。$K=4$ 那一幅显示同样
的换班,但多了一个轴:peak1 内出现可观比例的 $\texttt{C1pJ1p}$（用 shortcut +
jump-over）,peak2 内出现可观比例的 $\texttt{C0J1p}$（无 shortcut + jump-over)。
图~\ref{fig:mc-uncertainty-cn} 的 MC 置信区间确认这些组成差异显著大于 95\% CI,
所以这种反转是统计事实而非采样波动。

\textbf{机制。} 早峰是\emph{快通道}:walker 在出发不久便遭遇单向 shortcut,
近似直达 target,在小 $t$ 处吸收。晚峰是\emph{慢通道}:walker 始终未触发 shortcut,
只能沿环走"远路",其首达时间由未受扰动 ring 上的扩散时标决定。Valley 是两通道
质量都很少的时段——快通道已经吸收完毕、慢通道尚未开始大规模到达。对 $K=4$ 而言,
更宽的步长 $r=\lfloor K/2\rfloor=2$ 引入第二个分离轴:walker 可能\emph{跨过} target
而不被吸收(jump-over),从而推迟吸收、加重慢通道的尾部。这就是 $K=4$ 比 $K=2$
更稳定地呈现双峰的原因——它有两条独立的时标分离机制(shortcut 与 jump-over),
而 $K=2$ 只有一条。

\textbf{结论。} Lazy ring + shortcut 模型中 $f(t)$ 的双峰是\emph{两类首达时标
不同的轨迹群体的叠加},是结构性而非随机性的。peak1/peak2 窗口不是单纯的"曲线高处",
而是\emph{每一种群体占主导的参数区间};窗口内的类别比例正是这种解释的直接证据。

\section{敏感性分析（正文）}
\subsection{阈值敏感性}
阈值网格扫描覆盖 $h_{\min}\in\{10^{-8},10^{-7},10^{-6}\}$、
second\_frac$\in\{0.005,0.01,0.02\}$。
$t_2/t_1\in\{8,10,12\}$。\\
hv\_over\_max$\in\{0.10,0.20,0.30\}$。
图~\ref{fig:threshold-heatmap-cn} 用两类热图汇总通过率：
左侧为 hv\_over\_max 与 second\_frac 的二维热图（平均 $h_{\min},t_2/t_1$），\\
右侧为 $t_2/t_1$ 与 $h_{\min}$ 的二维热图（平均 hv\_over\_max 与 second\_frac）。
本图聚合 $N=100$ 且 $\beta$ 在 0.001--0.03 的 5 个取值：
$K=2$ 的平均通过率约 $0.33$，$K=4$ 约 $0.60$；
hv\_over\_max 越大通过率越高（$K=2$: 0.20/0.40/0.40；$K=4$: 0.40/0.60/0.80），
而 $h_{\min}$、second\_frac 与 $t_2/t_1$ 的影响较弱。
因此“$K=4$ 比 $K=2$ 更易双峰”的结论对阈值选择稳健，但对 hv\_over\_max 最敏感。

\begin{figure}[!htbp]
  \centering
  \includegraphics[width=0.8\linewidth]{outputs/sensitivity/threshold_sweep_heatmap.pdf}
  \caption{阈值敏感性热图（通过率）。}
  \label{fig:threshold-heatmap-cn}
\end{figure}

\subsection{窗口位置与宽度敏感性}
对每个窗口 $[t_L,t_R]$，分别进行整体平移与宽度变化：
$\Delta\in\{-5,-2,-1,0,1,2,5\}$（单位为离散时间步），
shift 模式为 $[t_L+\Delta, t_R+\Delta]$，width 模式为 $[t_L-\Delta, t_R+\Delta]$。
图~\ref{fig:bin-shift-cn} 同时展示两种模式。
当 cond\_by\_t 中存在缺失的时间点分类概率时，脚本使用对应窗口的基线组成进行插补，
以便评估边界扰动的影响。若窗口附近类别概率随 $t$ 近似单调，$\Delta$prop 会呈近似线性变化，
因此部分子图会出现“尖三角/折线”形态，这是窗口滑动的正常结果而非数值异常。
结果显示：$K=2$ 的最大比例变化 $\max|\Delta p|\le 4.49\times 10^{-3}$（平移）
与 $\le 4.26\times 10^{-3}$（宽度），
$K=4$ 的最大比例变化 $\max|\Delta p|\le 5.78\times 10^{-3}$（平移）
与 $\le 5.13\times 10^{-3}$（宽度）。
整体变化幅度仍小于 $1\%$ 量级，peak1/valley 中 $J1p$ 的稳定占比、peak2 由 $C0$ 类主导的结论保持不变。

\begin{figure}[!htbp]
  \centering
  \includegraphics[width=0.95\linewidth]{outputs/sensitivity/bin_shift_plot.pdf}
  \caption{窗口位置与宽度敏感性（shift/width 两种模式）。}
  \label{fig:bin-shift-cn}
\end{figure}

\subsection{MC 不确定性}
基于窗口内样本量，采用正态近似给出 95\% CI。
当 cond\_by\_t 无法覆盖窗口时，使用窗口级计数（summary/JSON）作为替代。
结果见图~\ref{fig:mc-uncertainty-cn}：$K=2$ 的 CI 平均宽度约 $0.011$，最大约 $0.036$，
$K=4$ 的 CI 平均宽度约 $0.029$，最大约 $0.063$。
这些区间仍显著小于类间差异，
不影响“$J1p$ 在 peak1/valley 中稳定存在”的判断。

\begin{figure}[!htbp]
  \centering
  \includegraphics[width=0.9\linewidth]{outputs/sensitivity/mc_uncertainty_plot.pdf}
  \caption{MC 不确定性（误差条/CI）。}
  \label{fig:mc-uncertainty-cn}
\end{figure}

\section*{小结}
Lazy ring + 单向 shortcut 模型中 $f(t)$ 的双峰是两类首达时标不同的轨迹群体的叠加:
\emph{快通道}(\texttt{C1p}\,$\cdot$,使用过 shortcut)贡献 $t_1$ 附近的早峰；
\emph{慢通道}(\texttt{C0}\,$\cdot$,从未使用 shortcut、沿环走远路)贡献 $t_2$ 附近的晚峰；
两峰之间的 valley 是两通道都不显著贡献质量的时段。图~\ref{fig:fig1-overlap-cn} 把这
个机制直接画了出来：peak1/peak2 的栈叠条在 $K=2$ 与 $K=4$ 都分别由 \texttt{C1p} /
\texttt{C0} 主导;$K=4$ 那一幅还多出一个分离轴(\texttt{J1p}, jump-over),
解释了它更稳定双峰的原因。

这一解读在三个维度上都稳健：(i)~阈值敏感性下,"$K=4$ 比 $K=2$ 更易双峰"的结论
跨越广泛的检测参数网格依然成立,只有 \texttt{hv\_over\_max} 引入不可忽略的变化;
(ii)~窗口位置与宽度扰动 $\Delta=\pm 1,2,5$ 仅引起 $<1\%$ 的比例变化,因此通道换班的
说法不依赖窗口边界的精确位置;(iii)~MC 置信区间在每个窗口内都明显小于类间差异,
所以"哪一类主导哪一峰"是统计上良好分离的事实。

模型推导、完整 $\beta$ 扫描以及更广 $K,N,q$ 范围请参阅 \texttt{reports/lazy\_jump\_ext/}。
姊妹报告 \texttt{ring\_lazy\_flux/} 通过 shortcut 通量给出另一种诊断,
\texttt{ring\_valley/} 则讨论无 shortcut 时的 valley 结构。

\end{document}
